\section{Operadores areal y modular en los sectores \texorpdfstring{\(90\)}{90} y \texorpdfstring{\(120\)}{120}}

En cada una de las dos familias de
\cref{rm:eq:modular-boundary-counts} existen dieciocho canales qutrit. Para
\(m\in\{90,120\}\), definimos
\begin{equation}
\label{rm:eq:modular-number-operators}
N_m=\sum_{e\in\partial A_m}N_e,
\qquad \Spec(N_e)=\{0,1,2\}.
\end{equation}
La salida de Barbero--Immirzi se usa aqu\'i como resultado previo, sin
rederivarla. Con la normalizaci\'on areal construida en
\cref{rm:eq:modular-theta-clock-a0}, resulta
\begin{equation}
\label{rm:eq:modular-area-operator}
\boxed{
\frac{\widehat{\mathcal A}_{\rm phys}}{\ell_P^2}
=\gammaH a_0(N_{90}+N_{120}),}
\qquad
\gammaH a_0=0.0016755940749224991.
\end{equation}

\begin{theorem}[Espectro finito del cuanto de \`area]
\label{rm:thm:modular-area-spectrum}
Sobre los treinta y seis enlaces del corte,
\begin{equation}
\label{rm:eq:modular-total-number-spectrum}
\Spec(N_{90}+N_{120})=\{0,1,\ldots,72\}.
\end{equation}
La multiplicidad del autovalor \(n\) es el coeficiente de \(z^n\) en
\((1+z+z^2)^{36}\). En consecuencia,
\begin{equation}
\label{rm:eq:modular-area-spectrum}
\boxed{
\Spec\!\left(\frac{\widehat{\mathcal A}_{\rm phys}}{\ell_P^2}\right)
=\{n\gammaH a_0:0\leq n\leq72\}.}
\end{equation}
Los cortes compatibles prolongan la familia
\(\mathcal A_n^{\rm prot}=n\gammaH a_0\ell_P^2\).
\end{theorem}

\begin{proof}
Cada operador \(N_e\) tiene funci\'on generatriz espectral
\(1+z+z^2\). Los treinta y seis factores conmutan y act\'uan sobre factores
tensoriales distintos; por tanto, la funci\'on generatriz del total es
\((1+z+z^2)^{36}\), cuyos grados recorren todos los enteros de cero a
setenta y dos. Multiplicar por el escalar positivo
\(\gammaH a_0\ell_P^2\) da \cref{rm:eq:modular-area-spectrum}.
\end{proof}

Es necesario distinguir dos procedimientos de cuantizaci\'on. La ley de corte
\cref{rm:thm:modular-triadic-cut} cuenta ochenta y una capacidades unitarias y
determina \(81\log3\); el operador
\cref{rm:eq:modular-area-operator} cuantiza ocupaciones de treinta y seis
enlaces de respuesta, hasta setenta y dos trits de n\'umero. Comparten la
misma hoja de borde, pero no el mismo observable.

La torsi\'on modular produce
\begin{equation}
\label{rm:eq:modular-delta-value}
\Dmod
=0.0001111970500599773249414566131933856\ldots,
\end{equation}
y el Hamiltoniano de borde
\begin{equation}
\label{rm:eq:modular-hamiltonian}
\boxed{
\HHM^{(A)}=-\log\rho_A
=cI+\Dmod(90N_{90}+120N_{120}).}
\end{equation}
El calificativo \emph{hologr\'afico} no designa una analog\'ia: la
incidencia \(6\to8\) se transporta a los conjuntos de borde
\(\partial A_{90}\) y \(\partial A_{120}\), y el estado reducido reside
en el \`algebra generada por sus operadores n\'umero.

\begin{theorem}[Reconstrucci\'on conjunta de geometr\'ia y memoria]
\label{rm:thm:modular-area-memory}
Sean
\begin{equation}
\label{rm:eq:modular-ND}
N=N_{90}+N_{120},
\qquad
D=90N_{90}+120N_{120}.
\end{equation}
Entonces
\begin{equation}
\label{rm:eq:modular-inverse-channels}
\boxed{
N_{90}=4N-\frac{D}{30},
\qquad
N_{120}=\frac{D}{30}-3N.}
\end{equation}
En particular,
\([\widehat{\mathcal A}_{\rm phys},\HHM^{(A)}]=0\), y el par
\((\widehat{\mathcal A}_{\rm phys},\HHM^{(A)}-cI)\) reconstruye los dos
sectores de ocupaci\'on.
\end{theorem}

\begin{proof}
Los operadores \(N_{90}\) y \(N_{120}\) act\'uan sobre factores distintos
y conmutan. El mapa lineal que lleva \((N_{90},N_{120})\) a \((N,D)\)
tiene matriz
\[
\begin{pmatrix}1&1\\90&120\end{pmatrix},
\qquad \det=30.
\]
Su inversa es exactamente \cref{rm:eq:modular-inverse-channels}. Como
\cref{rm:eq:modular-area-operator,rm:eq:modular-hamiltonian} son combinaciones
lineales de los mismos operadores n\'umero, el conmutador se anula.
\end{proof}

La interpretaci\'on operatoria es precisa: el observable de \`area conserva
el n\'umero total de excitaciones, mientras que el Hamiltoniano modular
conserva su sector de procedencia. Ninguno de ambos observables retiene por
s\'i solo las dos coordenadas; el par conjunto s\'i las determina.

\section{Confluencia de \`area, informaci\'on y energ\'ia}

La escala angular \(\theta_{\rm clk}\) de
\cref{rm:eq:modular-theta-clock-a0} y la temperatura cronol\'ogica
\(\Theta_{\rm clk}\) son objetos distintos; el uso de may\'uscula es parte
del tipado. La realizaci\'on termodin\'amica del \cref{rm:ch:metrology} demuestra
\begin{equation}
\label{rm:eq:modular-planck-trit-bridge}
E_P=k_B\Theta_{\rm clk}\log3,
\end{equation}
mientras la realizaci\'on areal construye el incremento elemental
\begin{equation}
\label{rm:eq:modular-area-quantum}
\delta\mathcal A
=\gammaH a_0\ell_P^2.
\end{equation}
No se igualan: \cref{rm:eq:modular-planck-trit-bridge} fija el cuanto
informacional de energ\'ia y \cref{rm:eq:modular-area-quantum} el cuanto
geom\'etrico de \`area. Su cociente define la tensi\'on de conversi\'on
tipada
\begin{equation}
\label{rm:eq:modular-information-area-tension}
\boxed{
\mathcal T_{I/A}
=\frac{E_P}{\delta\mathcal A}
=\frac{k_B\Theta_{\rm clk}\log3}
{\gammaH a_0\ell_P^2}.}
\end{equation}
La expresi\'on no introduce una constante adicional: determina las
normalizaciones que deben conservarse en la transformaci\'on de informaci\'on
en geometr\'ia.


En el corte m\'inimo de ochenta y un enlaces, la capacidad informacional
f\'isica y la energ\'ia de un trit por enlace son
\begin{equation}
\label{rm:eq:modular-cut-info-energy}
S_{\rm cut}^{\rm phys}=81k_B\log3,
\qquad
E_{\rm cut}^{(1)}=81E_P
=\Theta_{\rm clk}S_{\rm cut}^{\rm phys}.
\end{equation}
La igualdad final es una ecuaci\'on de estado de esta realizaci\'on, no una ley de
equipartici\'on universal. En particular, el estado Gibbs de los treinta y
seis enlaces del borde tiene entrop\'ia distinta y se trata a continuaci\'on.

Esta separaci\'on determina el contenido estructural de la confluencia:
Barbero--Immirzi
orienta el espectro de \`area; \(\log3\) mide la decisi\'on irreducible del
TRIT; \(\Theta_{\rm clk}\) convierte esa informaci\'on en energ\'ia; y
\(\ell_P^2\) conecta el cuanto de \`area con la familia gravitatoria. La
gravedad queda descrita como confluencia de realizaciones funcionales de un
antecedente com\'un, y no como una igualdad impuesta entre sus valores.
